%	These definitions were written by Simon Thompson, 1989-1991.
%
%	(c) Simon Thompson 1991.

% This file produced June 1994. Cut-down version of defs1.tex. Also
% includes new fonts and other changes.

\newfont{\latt}{cmtt12}
\newfont{\mytt}{cmtt10}
\newfont{\smtt}{cmtt8}
\newcommand{\mi}{\mytt}

% Force the separation between paragraphs to be zero.

\setlength{\parskip}{0pt}

%	First something to give Miranda source code listings.
%	Then a lot of abbreviations, to work in both LR and
%	math modes. These don't work in arrays, etc.
%	Finally, the basic macros for proof rules.

%	Definitions of environment, a change to the
%	obeyspaces macro (cf ref. for details) and
%	a number of definitions of new commands to
%	use quantifiers and logical operations outside
%	math mode.

%change of catcode commented out - 12.6.92
\newenvironment{ttdisplay}{\begin{alltt}%
\mi\obeylines\obeyspaces\frenchspacing\catcode`\_=12%
\setlength{\parskip}{0pt}%
}{\end{alltt}}

\newcommand{\so}{\begin{ttdisplay}\parindent 1pc}
\newcommand{\st}{\end{ttdisplay} \noindent }
\newcommand{\stt}{\end{ttdisplay}}

\newcommand{\inso}{\mi\catcode`\_=12}
\newcommand{\inst}{\rm}

% An environment which allows Miranda text to be printed in a verbatim
% form, whilst at the same time allowing math style quantifiers and
% the like.
% The change of catcode for underscore is to disable its
% special interpretation under TeX. We can therefore
% still use it within Miranda identifiers, etc.
% The abbreviations \so \st are used for historical reasons, as
% similar ones were used in troff -ms.
% The \noindent in the \st is to ensure that we can use \st
% within a paragraph to perform some displaying, without
% the line that follows starting with an indentation.

{\obeyspaces\global\let =\ }

% See pages 380-381 of the TeX book for the source of the line above.
% It makes sure that spaces are really treated as spaces by the 
% obeyspaces macro, and so by the ttdisplay environment.

% Tight itemized lists: make \itemsep zero

\newenvironment{titemize}{\begin{list}
                        {$\bullet$}
                         {\setlength{\itemsep}{0pt}}}{\end{list}}

% Logical symbols -- uses the 2e \ensuremath command

\newcommand{\all}{\ensuremath{\forall}}
\newcommand{\allv}{\ensuremath{\forall_{\hspace*{-1.5pt}v\hspace*{1pt}}}}
\newcommand{\exi}{\ensuremath{\exists}}
\newcommand{\ou}{\ifmmode \vee \else$\vee$\fi}
\newcommand{\an}{\ensuremath{\wedge}}
\newcommand{\imp}{\ensuremath{\Rightarrow}}
\newcommand{\no}{\ensuremath{\neg}}
\newcommand{\bi}{\ensuremath{\Leftrightarrow}}
\newcommand{\turn}{\ensuremath{\vdash}}
\newcommand{\bo}{\ensuremath{\bot}}
\newcommand{\se}{\ensuremath{\equiv}}

% Logical symbols for 209: commented out

%\newcommand{\an}{\ifmmode \wedge \else$\wedge$\fi}
%\newcommand{\all}{\ifmmode \forall \else$\forall$\fi}
%\newcommand{\exi}{\ifmmode \exists \else$\exists$\fi}
%\newcommand{\imp}{\ifmmode \Rightarrow \else$\Rightarrow$\fi}
%\newcommand{\no}{\ifmmode \neg \else$\neg$\fi}
%\newcommand{\bi}{\ifmmode \Leftrightarrow \else$\Leftrightarrow$\fi}
%\newcommand{\turn}{\ifmmode \vdash \else$\vdash$\fi}
%\newcommand{\bo}{\ifmmode \bot \else$\bot$\fi}
%\newcommand{\se}{\ifmmode \equiv \else$\equiv$\fi}

% Backslash and Or in mi font

\newcommand{\bs}[1]{{\ttfamily \symbol{92}#1}}
\newcommand{\mor}{{\ttfamily \symbol{92}/}}

% Miranda quantifiers
 
%\newcommand{\alldf}{\ifmmode \forall_{\tt\,\!def} \else $\forall_{\tt\,\!def}$\fi}
%\newcommand{\allfdf}{\ifmmode \forall_{\tt fin-def} \else $\forall_{\tt fin-def}$\fi}
%\newcommand{\allfin}{\ifmmode \forall_{\tt fin} \else $\forall_{\tt fin}$\fi}

% A proper twiddle, with height adjusted.

\newcommand{\twid}{\raisebox{0ex}[0ex][0ex]{\Large\raisebox{-5pt}{$\,\tilde{}\:$}}}

% Comment symbol in a draft.

\newcommand{\comment}[1]{\noindent $\bigcirc$ {\large \sl #1}}

% Exercises

\newcommand{\exer}{\vspace{0.1in}\noindent{\textsf{\textbf{
  Exercises}}}\nopagebreak}%\vspace{0.1in}\nopagebreak}
\newcommand{\exone}{\vspace{0.1in}\noindent{\textsf{\textbf{
  Exercise}}}
 \nopagebreak\vspace{0.1in}\nopagebreak}
\newcommand{\case}{\medskip\noindent{\bf Case}}

% Numbering Exercises. Restart numbering in each chapter
% -- this needs to be done by a restart at the beginning
% of each chapter. See \chapstart below

% \exn is the command to give the next exercise number

\newcounter{exnum}

\newcommand{\exn}{\medskip\noindent\addtocounter{exnum}{1}\arabic{chapter}.\arabic{exnum}\ }

% Cross-referencing of exercises. NB only allows backwards references!

% \sexn{XXX} generates the next exercise number and saves the typeset text
% in the \sbox called XXX; available boxes currently are \locala, \localb
% \sexn uses \exnshort, not used elsewhere.

\newcommand{\exnshort}{\addtocounter{exnum}{1}\arabic{chapter}.\arabic{exnum}}
\newcommand{\sexn}[1]{\sbox{#1}{\exnshort}\medskip\noindent\usebox{#1}}

% Chapter start macro - reset exercise counter to zero.

\newcommand{\chapstart}{\setcounter{exnum}{0}}

% Tables

\newenvironment{mytable}{\begin{tabular}{>{\hskip1pc}p{1.3in}p{3in}}}{\end{tabular}}
\newenvironment{mytablethree}{\begin{tabular}{p{0.7in}p{1.5in}p{2in}}}{\end{tabular}}

% Local references: use the boxes \locala \localb

\newsavebox{\locala}
\newsavebox{\localb}

% Drawing a vertical line: the argument is its height. It has zero width
% and is defined not to extend beyond the line on which it lies.

\newcommand{\vertline}[1]{
\hspace*{2pt}\raisebox{0pt}[0pt][0pt]{\rule{0.5pt}{#1}}\hspace*{-2.5pt}} 

% start and end of proof

\newcommand{\startpr}{\subparagraph{Proof}}
\newcommand{\blackbox}{\hfill\rule{8pt}{8pt}}  

% explaining an equation

\newcommand{\push}[1]{\hfill{\rm #1}}


% Natural deduction rule layout.

\newcommand{\tworule}[3]
{\begin{array}{cc}
#1 & #2 \\
\hline
\multicolumn{2}{c}{#3}
\end{array} }

\newcommand{\threerule}[4]
{\begin{array}{ccc}
#1 & #2 & #3 \\
\hline
\multicolumn{3}{c}{#4}
\end{array} }


% Semantic brackets (1/98)

\newcommand{\seml}{\ensuremath{[\![}}
\newcommand{\semr}{\ensuremath{]\!]}}

% One step of a calculation (6/98)

\newcommand{\step}{{\ensuremath\leadsto\ }}

